\documentclass[../../main.tex]{subfiles}
\graphicspath{{../../image/}{../../figure/}}
\begin{document}

\section{周期解}\label{section:ch12-周期解}

真非线性两守恒律组的研究将以对具小振荡的周期解大时间行为的讨论作结：
\begin{align}
U(x+\ell,t)=U(x,t),\qquad -\infty<x<\infty,\ t>0,
\label{align:ch12-周期}
\end{align}
且零均值\footnote{若 Cauchy 问题有唯一解，具零均值的周期初值必然生成具有同一性质的解。}：
\begin{align}
\int_y^{y+\ell}U(x,t)\,dx=0,\qquad -\infty<y<\infty,\ t>0.
\label{align:ch12-零均值}
\end{align}

周期性导致的波禁闭诱导主动相互作用与抵消。结果，每周总变差以 $O(t^{-1})$ 的速率衰减：

\begin{theorem}\label{theorem:ch12-周期TV}
对任意 $x\in(\infty,\infty)$ 及 $t>0$，
\begin{align}
TV_{[x,x+\ell]}z(\cdot,t)+TV_{[x,x+\ell]}w(\cdot,t)\leqslant\frac{b\ell}{t}.
\label{align:ch12-周期TV}
\end{align}
\end{theorem}

\begin{proof}
对 $y=x+n\ell$ 应用\eqref{align:ch12-TV散开}，然后除以 $n$ 并令 $n\to\infty$。证明完成。
\end{proof}

现我们在 $O(t^{-1})$ 尺度上求解渐近性。在原书第 11.7 节真非线性标量守恒律语境中遇到的机制，
即极值后向特征线在周期长度的 $x$ 轴区间上的交点禁闭，在此同样有效并生成类似的锯齿形渐近剖面。剖面的节点再次被分线（divide）跟踪，即定义 10.3.3 意义下的分线。

\begin{theorem}\label{theorem:ch12-周期渐近}
上半平面被极小（或极大）1-（或 2-）分线划分，沿之分线 $z$（或 $w$）当 $t\to\infty$ 时快速衰减为零，为 $O(t^{-2})$。设 $\chi_-(\cdot)$ 与 $\chi_+(\cdot)$ 是任意两条相邻 1-（或 2-）分线，$\chi_-(t)<\chi_+(t)$。则 $\chi_+(t)-\chi_-(t)$ 当 $t\to\infty$ 时以 $O(t^{-1})$ 的速率逼近常数。进而，在 $\chi_-$ 与 $\chi_+$ 之间存在一条 1-（或 2-）特征线 $\psi$，使当 $t\to\infty$ 时，
\begin{align}
\psi(t)=\tfrac12[\chi_-(t)+\chi_+(t)]+o(1),
\label{align:ch12-中线}
\end{align}
\begin{align}
\lambda_z(0,0)z(x,t)=
\begin{cases}
\dfrac{x-\chi_-(t)}{t}+o\left(\dfrac1t\right),\qquad &\chi_-(t)<x<\psi(t),\\[2mm]
\dfrac{x-\chi_+(t)}{t}+o\left(\dfrac1t\right),\qquad &\psi(t)<x<\chi_+(t),
\end{cases}
\label{align:ch12-周期z}
\end{align}
或
\begin{align}
\mu_w(0,0)w(x,t)=
\begin{cases}
\dfrac{x-\chi_-(t)}{t}+o\left(\dfrac1t\right),\qquad &\chi_-(t)<x<\psi(t),\\[2mm]
\dfrac{x-\chi_+(t)}{t}+o\left(\dfrac1t\right),\qquad &\psi(t)<x<\chi_+(t).
\end{cases}
\label{align:ch12-周期w}
\end{align}
\end{theorem}

证明上述命题的第一步是考察分线的大时间行为：

\begin{lemma}\label{lemma:ch12-分线衰减}
沿极小（或极大）1-（或 2-）分线，$z$（或 $w$）当 $t\to\infty$ 时以 $O(t^{-2})$ 的速率衰减。进而，若 $\chi_-(\cdot)$ 与 $\chi_+(\cdot)$ 是任意两条极小（或极大）1-（或 2-）分线，则当 $t\to\infty$ 时，
\begin{align}
\chi_+(t)-\chi_-(t)=h_\infty+O\left(\frac1t\right),
\label{align:ch12-分线间距}\\
\int_{\chi_-(t)}^{\chi_+(t)}z(x,t)\,dx=O\left(\frac1{t^2}\right),
\label{align:ch12-分线积分1}
\end{align}
或
\begin{align}
\int_{\chi_-(t)}^{\chi_+(t)}w(x,t)\,dx=O\left(\frac1{t^2}\right).
\label{align:ch12-分线积分2}
\end{align}
\end{lemma}

\begin{proof}
设 $\chi(\cdot)$ 是极小 1-分线，即从点 $\{(x_n,t_n)\}$（$t_n\to\infty$，当 $n\to\infty$）发出的极小后向 1-特征线序列 $\{\xi_n(\cdot)\}$ 的极限。令 $(z_n(\cdot),w_n(\cdot))$ 表示沿 $\xi_n(\cdot)$ 的 $(z,w)$ 迹。将\cref{theorem:ch12-TV界}用于时间原点平移到 $\tau$ 的情形，并利用\cref{theorem:ch12-周期TV}，推得 $w_n(\cdot)$ 在任一区间 $[\tau,\tau+1]\subset[0,t_n]$ 上的总变差为 $O(\tau^{-1})$，关于 $n$ 一致。因此，由\cref{theorem:ch12-传播主定理}，$z_n(\cdot)$ 是 $[0,t_n]$ 上的非增函数，其在 $[\tau,\tau+1]$ 上的振荡为 $O(\tau^{-3})$，关于 $n$ 一致。由此，沿 $\chi(\cdot)$ 的 $z$ 迹 $\bar z(\cdot)$ 同样是非增函数，在 $[\tau,\tau+1]$ 上振荡为 $O(\tau^{-3})$。通过将 $\bar z(\cdot)$ 的振荡按单位长度区间从 $t$ 到无穷总计，我们验证断言 $\bar z(t)=O(t^{-2})$。

类似论证表明沿极大 2-分线的 $w$ 迹 $\bar w(\cdot)$ 当 $t\to\infty$ 时同样为 $O(t^{-2})$。

设 $\chi_-(\cdot)$ 与 $\chi_+(\cdot)$ 是极小 1-分线，$h(t)=\chi_+(t)-\chi_-(t)\geqslant0$（$0\leqslant t<\infty$）。注意，由于周期性，若 $h(0)<k\ell$（某整数 $k$），则 $h(t)\leqslant k\ell$（$0\leqslant t<\infty$）。令 $(z_-(\cdot),w_-(\cdot))$ 与 $(z_+(\cdot),w_+(\cdot))$ 分别表示沿 $\chi_-(\cdot)$ 与 $\chi_+(\cdot)$ 的 $(z,w)$ 迹，我们有
\begin{align}
\dot h(\tau)=\lambda(z_+(\tau),w_+(\tau))-\lambda(z_-(\tau),w_-(\tau)),
\label{align:ch12-hderiv}
\end{align}
对几乎所有 $\tau\in[0,\infty)$ 成立。

从点 $(\chi_+(\tau),\tau)$ 发出的极大后向 2-特征线 $\zeta_\tau(\cdot)$ 在时刻 $\tau-f(\tau)$ 被 $\chi_-(\cdot)$ 的图像截取。若 $(\hat z(\cdot),\hat w(\cdot))$ 表示沿 $\zeta_\tau(\cdot)$ 的 $(z,w)$ 迹，则\cref{theorem:ch12-TV界}与\cref{theorem:ch12-周期TV}共同蕴含 $\hat z(\cdot)$ 在区间 $[\tau-f(\tau),\tau]$ 上的总变差为 $O(\tau^{-1})$（当 $\tau\to\infty$）。于是由\cref{theorem:ch12-传播主定理}，$\hat w(\cdot)$ 在 $[\tau-f(\tau),\tau]$ 上的振荡为 $O(\tau^{-3})$。因此
\begin{align}
w_+(\tau)=w_-(\tau-f(\tau))+O(\tau^{-3}).
\label{align:ch12-w+}
\end{align}
由于 $z_\pm(\tau)=O(\tau^{-2})$，\eqref{align:ch12-hderiv}给出
\begin{align}
\dot h(\tau)=\lambda(0,w_-(\tau-f(\tau)))-\lambda(0,w_-(\tau))+O\left(\frac1{\tau^2}\right).
\label{align:ch12-hderiv2}
\end{align}

由 $\dot h(\tau)=O(\tau^{-1})$ 与 $\dot\zeta_\tau=\mu(0,0)+O(\tau^{-1})$，推得 $f(\cdot)$ 在区间 $[\tau,\tau+1]$ 上的振荡为 $O(\tau^{-1})$。$w_-(\cdot)$ 在 $[\tau,\tau+1]$ 上的总变差同样为 $O(\tau^{-1})$。于是对任意 $t<t'<\infty$，
\begin{align}
\left|\int_t^{t'}\{\lambda(0,w_-(\tau-f(\tau)))-\lambda(0,w_-(\tau))\}\,d\tau\right|\leqslant\frac{c}{t}.
\label{align:ch12-h积分}
\end{align}
结合\eqref{align:ch12-hderiv2}与\eqref{align:ch12-h积分}，即得\eqref{align:ch12-分线间距}。

令 $U_-(\cdot)$ 与 $U_+(\cdot)$ 分别表示沿 $\chi_-(\cdot)$ 与 $\chi_+(\cdot)$ 的 $U$ 迹。将\eqref{align:ch12-系统}在 $\{(x,\tau):t<\tau<\infty,\ \chi_-(\tau)<x<\chi_+(\tau)\}$ 上积分，
得方程
\begin{align}
\int_{\chi_-(t)}^{\chi_+(t)}U(x,t)\,dx
=\int_t^\infty\{F(U_+(\tau))-\lambda(U_+(\tau))U_+(\tau)-F(U_-(\tau))+\lambda(U_-(\tau))U_-(\tau)\}\,d\tau.
\label{align:ch12-U积分}
\end{align}

我们用行向量 $\mathrm{D}z(0)$ 左乘\eqref{align:ch12-U积分}。由\eqref{align:ch12-riemann归一}，
$U_z=R$、$U_w=S$，故利用\eqref{align:ch12-速矢}推得
\begin{align}
\mathrm{D}z(0)U&=z+O(z^2+w^2),
\label{align:ch12-DzU}\\
\mathrm{D}z(0)[F(U)-\lambda(U)U]&=\mathrm{D}z(0)F(0)+aw^2+O(z^2+|zw|+|w|^3),
\label{align:ch12-DzF}
\end{align}
其中常数 $a$ 是 $\frac12(\lambda-\mu)S^\top\mathrm{D}^2_{z}S$ 在 $U=0$ 处的值。结合 $z_\pm(\tau)=O(\tau^{-2})$、$w_\pm(\tau)=O(\tau^{-1})$ 与\eqref{align:ch12-w+}，我们断定
\begin{align}
\int_{\chi_-(t)}^{\chi_+(t)}z(x,t)\,dx=a\int_t^\infty[w_-^2(\tau-f(\tau))-w_-^2(\tau)]\,d\tau+O\left(\frac1{t^2}\right).
\label{align:ch12-z积分-s10}
\end{align}
如上所释，在区间 $[\tau,\tau+1]$ 上 $f(\cdot)$ 的振荡为 $O(\tau^{-1})$，$w_-^2(\cdot)$ 的总变差为 $O(\tau^{-2})$。于是\eqref{align:ch12-z积分}右端当 $t\to\infty$ 时为 $O(t^{-2})$，这就建立了\eqref{align:ch12-分线积分1}。

当 $\chi_-(\cdot)$ 与 $\chi_+(\cdot)$ 是极大 2-分线时，类似论证验证\eqref{align:ch12-分线间距}与\eqref{align:ch12-分线积分2}。证明完成。
\end{proof}

\cref{theorem:ch12-周期渐近}的其余断言将通过下述引理建立：

\begin{lemma}\label{lemma:ch12-中线}
考虑任意两条相邻极小（或极大）1-（或 2-）分线 $\chi_-(\cdot)$、$\chi_+(\cdot)$，$\chi_-(t)<\chi_+(t)$（$0\leqslant t<\infty$）。在\cref{section:ch12-解的局部结构}的记号中，从任一固定点 $(\bar x,0)$（$\chi_-(0)<\bar x<\chi_+(0)$）发出的特殊前向 1-（或 2-）特征线 $\phi_-(\cdot)$（或 $\psi_+(\cdot)$）记为 $\psi(\cdot)$。则 $\psi(\cdot)$ 满足\eqref{align:ch12-中线}。进而\eqref{align:ch12-周期z}（或\eqref{align:ch12-周期w}）成立。
\end{lemma}

\begin{proof}
只需讨论 $\chi_-$ 与 $\chi_+$ 为 1-分线的情形。我们考虑从点 $(x,t)$（$t>0$，$\chi_-(t)<x<\chi_+(t)$）发出的极小后向 1-特征线 $\xi(\cdot)$。它们的图像被限制在 $\chi_-$ 与 $\chi_+$ 的图像之间。此类 $\xi$ 与 $x$ 轴的交点 $\xi(0)$ 不能在开区间 $(\chi_-(0),\chi_+(0))$ 内聚到任一点 $\hat x$，因为在此情形将从点 $(\hat x,0)$ 发出一条极小 1-分线，与 $\chi_-$、$\chi_+$ 相邻的假设矛盾。因此，由 $\psi(\cdot)$ 的构造推知，当 $t\to\infty$ 时，若 $x\in(\chi_-(t),\psi(t)]$ 则 $\xi(t)\to\chi_-(\tau)$；若 $x\in(\psi(t),\chi_+(t)]$ 则 $\xi(t)\to\chi_+(\tau)$，收敛在 $[0,\infty)$ 的紧子集上一致。

现固定从某点 $(x,t)$（$\chi_-(t)<x\leqslant\psi(t)$）发出的 $\xi(\cdot)$，令 $h(\tau)=\xi(\tau)-\chi_-(\tau)$（$0\leqslant\tau\leqslant t$）。于是对几乎所有 $\tau\in[0,t]$，我们有
\begin{align}
\dot h(\tau)=\lambda(\bar z(\tau),\bar w(\tau))-\lambda(z_-(\tau),w_-(\tau)),
\label{align:ch12-hderiv3}
\end{align}
其中 $(\bar z(\cdot),\bar w(\cdot))$ 表示沿 $\xi(\cdot)$ 的 $(z,w)$ 迹，$(z_-(\cdot),w_-(\cdot))$ 表示沿 $\chi_-(\cdot)$ 的 $(z,w)$ 迹。

由\cref{theorem:ch12-TV界}与\cref{theorem:ch12-周期TV}，$\bar w(\cdot)$ 在任一区间 $[s,s+1]\subset[0,t]$ 上的总变差为 $O(s^{-1})$。于是由\cref{theorem:ch12-传播主定理}，$\bar z(\cdot)$ 在 $[s,s+1]$ 上的振荡为 $O(s^{-3})$，故
\begin{align}
\bar z(\tau)=z(x,t)+O\left(\frac1{\tau^2}\right).
\label{align:ch12-zbar}
\end{align}
进而，由\cref{lemma:ch12-分线衰减}，$z_-(\tau)=O(\tau^{-2})$。又 $z(x,t)=O(t^{-1})$，故 a fortiori $z(x,t)=O(\tau^{-1})$。由这些观察，\eqref{align:ch12-hderiv3}给出
\begin{align}
\dot h(\tau)=\lambda_z(0,0)z(x,t)+\lambda(0,\bar w(\tau))-\lambda(0,w_-(\tau))+O\left(\frac1{\tau^2}\right).
\label{align:ch12-hderiv4}
\end{align}

对任意固定 $\tau\gg0$，我们考虑从点 $(\xi(\tau),\tau)$ 发出的极大后向 2-特征线 $\zeta_\tau(\cdot)$，它在时刻 $\tau-f(\tau)$ 被 $\chi_-(\cdot)$ 的图像截取。若 $(\hat z(\cdot),\hat w(\cdot))$ 表示沿 $\zeta_\tau(\cdot)$ 的 $(z,w)$ 迹，则\cref{theorem:ch12-TV界}与\cref{theorem:ch12-周期TV}共同蕴含 $\hat z(\cdot)$ 在区间 $[\tau-f(\tau),\tau]$ 上的总变差为 $O(\tau^{-1})$。于是由\cref{theorem:ch12-传播主定理}，$\hat w(\cdot)$ 在 $[\tau-f(\tau),\tau]$ 上的振荡为 $O(\tau^{-3})$。因此
\begin{align}
\bar w(\tau)=w_-(\tau-f(\tau))+O\left(\frac1{\tau^3}\right),
\label{align:ch12-wbar}
\end{align}
故\eqref{align:ch12-hderiv4}蕴含
\begin{align}
\dot h(\tau)=\lambda_z(0,0)z(x,t)+\lambda(0,w_-(\tau-f(\tau)))-\lambda(0,w_-(\tau))+O\left(\frac1{\tau^2}\right).
\label{align:ch12-hderiv5}
\end{align}

如\cref{lemma:ch12-分线衰减}的证明，在任一区间 $[\tau,\tau+1]\subset[0,t]$ 上 $f(\cdot)$ 的振荡为 $O(\tau^{-1})$，$w_-(\cdot)$ 的总变差也为 $O(\tau^{-1})$。因此，将\eqref{align:ch12-hderiv5}在区间 $[s,t]$（$0<s<t$）上积分，推得
\begin{align}
x-\chi_-(t)-\lambda_z(0,0)z(x,t)t=\xi(s)-\chi_-(s)+O\left(\frac1s\right)+sO\left(\frac1t\right).
\label{align:ch12-h积分2}
\end{align}

关于\eqref{align:ch12-h积分2}右端，给定 $\varepsilon>0$，先取 $s$ 充分大使 $O(s^{-1})<\frac13\varepsilon$。$s$ 如此取定后，确定 $\hat t$，使当 $t\geqslant\hat t$ 时 $sO(t^{-1})$ 不超过 $\frac13\varepsilon$，同时对所有 $x\in(\chi_-(t),\psi(t)]$ 有 $\xi(s)-\chi_-(s)<\frac13\varepsilon$。只需对 $t=\hat t$、$x=\psi(\hat t)$ 检验后一条件。这样我们就验证了\eqref{align:ch12-h积分2}左端当 $t\to\infty$ 时在 $(\chi_-(t),\psi(t))$ 上关于 $x$ 一致为 $o(1)$，这就验证了\eqref{align:ch12-周期z}的上半部分。\eqref{align:ch12-周期z}的下半部分由类似论证建立。证明完成。
\end{proof}

\end{document}
